% chapters/requirements/rnf.tex
% Requisitos no funcionales
% ============================================================================

% RNF-1 --------------------------------------------------------------------
\item \textbf{Rendimiento.} \label{rnf:performance}
El sistema debe satisfacer los siguientes tiempos de respuesta en condiciones normales de
operación: las operaciones \acs{crud} estándar de la \acs{api} deben completarse en menos
de 500~ms; la carga y descarga de ficheros \acs{pdf} debe completarse en menos de
3~segundos; los listados paginados deben servirse en menos de 1~segundo; las exportaciones
de datos deben completarse en menos de 10~segundos; y las tareas asíncronas de \acs{ocr}
deben procesarse en menos de 60~segundos por instancia. Estos tiempos se miden desde la
recepción de la petición hasta el envío de la respuesta, excluyendo la latencia de red.

% RNF-2 --------------------------------------------------------------------
\item \textbf{Concurrencia.} \label{rnf:concurrency}
El sistema debe soportar un pico de concurrencia de aproximadamente 1\,200~usuarios
simultáneos por organización, correspondiente al escenario de publicación de notas o
apertura de revisión con cuatro asignaturas de 300~alumnos cada una. En operación normal de
corrección, el sistema debe soportar al menos 30~correctores simultáneos por organización
sin degradación apreciable del rendimiento. En modo \dfn{multi-tenant}, el sistema debe
aislar el rendimiento entre organizaciones para que un pico en una no afecte a las demás.

% RNF-3 --------------------------------------------------------------------
\item \textbf{Despliegue.} \label{rnf:deployment}
El sistema debe poder desplegarse mediante un único fichero \texttt{docker-compose.yml} que
levante todos los servicios necesarios: Django (\acs{api}), PostgreSQL (base de datos),
Redis (caché y \textit{broker} de colas), uno o más \dfnpl{worker} de Celery (\acs{ocr} y
notificaciones), MinIO (almacenamiento de objetos) y un servidor \acs{sftp} (ingesta de
\acs{pdf}). La configuración específica del entorno (credenciales, claves, puertos) debe
externalizarse en ficheros de variables de entorno. El sistema debe poder ejecutarse
íntegramente en una única máquina para el modo instalación institucional. En la versión~1.0,
una única instancia de Django es suficiente para este modo.

% RNF-4 --------------------------------------------------------------------
\item \textbf{Escalabilidad.} \label{rnf:scalability}
La arquitectura de contenedores debe facilitar un futuro despliegue distribuido en servicios
\textit{cloud} para el modo SaaS, sin cambios arquitectónicos en el código fuente. Los
\dfnpl{worker} de Celery deben poder replicarse horizontalmente. El aislamiento entre
organizaciones se garantiza a nivel lógico mediante el campo \texttt{organization\_id} y un
\textit{middleware} de Django que lo inyecta en las consultas.

% RNF-5 --------------------------------------------------------------------
\item \textbf{Disponibilidad y observabilidad.} \label{rnf:availability}
El sistema debe exponer un \textit{endpoint} de \textit{health check} (\texttt{/health/})
que verifique la conectividad con PostgreSQL, Redis, MinIO y el estado de los
\dfnpl{worker} de Celery. La respuesta debe indicar el estado individual de cada
componente y un estado global del sistema. Este \textit{endpoint} no requiere
autenticación.

% RNF-6 --------------------------------------------------------------------
\item \textbf{Seguridad en la \acs{api}.} \label{rnf:api_security}
El sistema debe implementar \textit{rate limiting} en la \acs{api} para prevenir abuso,
utilizando el mecanismo de \textit{throttling} de \acs{drf} o una solución equivalente.
Los límites deben ser configurables por organización y por tipo de usuario. Todas las
entradas del usuario deben validarse y sanearse para prevenir inyecciones y ataques
comunes. Las respuestas que contengan datos sensibles (\acs{pdf}, calificaciones, datos
personales descifrados) deben incluir cabeceras \acs{http} anti-caché
(\texttt{Cache-Control: no-store, no-cache, must-revalidate}, \texttt{Pragma: no-cache},
\texttt{Expires: 0}) para instruir al navegador y a cualquier \textit{proxy} intermedio a
no almacenar el contenido. En futuras versiones se implementará cifrado en tránsito
mediante HTTPS/TLS en todas las comunicaciones.

% RNF-7 --------------------------------------------------------------------
\item \textbf{Cumplimiento del \acs{rgpd}.} \label{rnf:gdpr}
El sistema debe cumplir con los requisitos técnicos derivados del \acs{rgpd} aplicables a
los datos que maneja (\acs{dni}, calificaciones, datos de posibles menores): cifrado en
reposo de datos especialmente sensibles (\acs{dni}), auditoría de accesos tanto a
modificaciones como a lecturas de datos de categoría especial, no eliminación física de
usuarios (\textit{soft delete}) y almacenamiento seguro de contraseñas. La anonimización
de registros y los \textit{endpoints} de gestión de derechos ARCO se contemplan para
versiones futuras. En el escenario de autodespliegue, la responsabilidad del tratamiento
recae sobre la institución; en el escenario SaaS futuro, se requerirán contratos de encargo
de tratamiento y despliegue en regiones del Espacio Económico Europeo.

% RNF-8 --------------------------------------------------------------------
\item \textbf{Interoperabilidad.} \label{rnf:interop}
El \dfn{backend} debe exponer una \acs{api} \acs{rest} versionada mediante prefijo de ruta
(\texttt{/api/v1/}) cuyo contrato se documenta mediante una especificación OpenAPI~3.x
generada automáticamente a partir del código fuente (mediante \texttt{drf-spectacular} o
herramienta equivalente). El sistema debe exponer la especificación en \textit{endpoints}
de documentación interactiva (Swagger~UI o ReDoc). Todas las respuestas de la \acs{api}
utilizan formato \acs{json}, salvo las descargas de ficheros binarios. Las respuestas de
error siguen un formato uniforme con código de error simbólico, sin mensajes en lenguaje
natural, para facilitar la independencia lingüística. El contrato \acs{api} constituye la
interfaz formal entre el \dfn{backend} y el \dfn{frontend}.

% RNF-9 --------------------------------------------------------------------
\item \textbf{Mantenibilidad y extensibilidad.} \label{rnf:maintainability}
El sistema debe diseñarse con una arquitectura modular que facilite la incorporación de
nuevas funcionalidades. En concreto, el autenticador y el motor de \acs{ocr} deben
implementarse como \dfnpl{plugin} intercambiables. El código debe seguir los principios de
separación de responsabilidades, con capas diferenciadas para la \acs{api} (serializadores
y vistas), la lógica de negocio (servicios) y el acceso a datos (modelos y repositorios).
Se debe evitar el acoplamiento directo entre subsistemas, favoreciendo la comunicación
mediante eventos o interfaces definidas. El filtrado \dfn{multi-tenant} por
\texttt{organization\_id} debe implementarse mediante un \textit{middleware} de Django que
inyecte automáticamente el filtro en todas las consultas.

% RNF-10 -------------------------------------------------------------------
\item \textbf{Internacionalización.} \label{rnf:i18n}
La interfaz de la \acs{api} debe ser lo más agnóstica al idioma posible: los códigos de
error, tipos de evento y claves de configuración se expresan en inglés técnico. Los
contenidos generados para el usuario (correos electrónicos, mensajes de notificación) deben
soportar internacionalización (i18n) mediante el sistema de traducción de Django. En la
versión~1.0, el idioma por defecto por organización es el español, con posibilidad de añadir
idiomas adicionales en versiones futuras mediante ficheros de traducción. El motor de
\acs{ocr} debe poder configurarse con el idioma del texto a reconocer.

% RNF-11 -------------------------------------------------------------------
\item \textbf{Observabilidad y \textit{logging}.} \label{rnf:logging}
El sistema debe mantener dos niveles de \textit{logging} diferenciados. El \textit{log} de
aplicación registra eventos técnicos (errores, advertencias, información de depuración)
siguiendo las prácticas estándar de cualquier aplicación en producción, con formato
estructurado (\acs{json}) y niveles de severidad configurables. El \textit{log} de
auditoría registra eventos de negocio con impacto en datos (detallado en
\cref{rf:sec:audit}). Ambos \textit{logs} son independientes y tienen políticas de
retención distintas: el \textit{log} de aplicación se rota periódicamente; el \textit{log}
de auditoría es inmutable y permanente.

% RNF-12 -------------------------------------------------------------------
\item \textbf{Pruebas.} \label{rnf:testing}
El sistema debe contar con una cobertura de pruebas que garantice la fiabilidad del
\textit{software}. Se requieren tres niveles de pruebas: pruebas unitarias para la lógica
de negocio y funciones de utilidad; pruebas de integración para verificar la interacción
entre componentes (base de datos, cola de tareas, almacenamiento de objetos, \acs{api},
filtrado \dfn{multi-tenant}); y pruebas de carga para validar los requisitos de rendimiento
y concurrencia bajo los escenarios definidos (picos de publicación, ingesta masiva de
\acs{pdf}).

% RNF-13 -------------------------------------------------------------------
\item \textbf{Persistencia e integridad de datos.} \label{rnf:persistence}
La base de datos principal es PostgreSQL. El sistema debe garantizar la integridad
referencial de los datos en todo momento: los usuarios no se eliminan físicamente
(\textit{soft delete} con \texttt{is\_active}), las relaciones entre entidades se mantienen
mediante claves foráneas con restricciones adecuadas y las operaciones críticas se ejecutan
dentro de transacciones atómicas. Redis se utiliza como caché, \textit{broker} de Celery y
almacén de la lista negra de \textit{tokens}; no se emplea como almacenamiento primario de
datos de negocio. El aislamiento \dfn{multi-tenant} se garantiza mediante el campo
\texttt{organization\_id} en todas las tablas principales y un \textit{middleware} que lo
inyecta en las consultas.

% RNF-14 -------------------------------------------------------------------
\item \textbf{Almacenamiento de objetos.} \label{rnf:object_storage}
MinIO se utiliza como servicio de almacenamiento de objetos para todos los ficheros de
tamaño variable: \acs{pdf} de instancias, \acs{pdf} en blanco de modelos, audios de
anotaciones y, en el futuro, imágenes. El almacenamiento en MinIO es independiente de la
base de datos: los metadatos de los ficheros (tamaño, tipo, referencia) se almacenan en
PostgreSQL y los ficheros binarios en MinIO. La organización de \dfnpl{bucket} debe
permitir el aislamiento \dfn{multi-tenant} (un \dfn{bucket} por organización o un prefijo
por organización dentro de un \dfn{bucket} compartido). Esta separación facilita el
escalado independiente del almacenamiento y la migración futura a servicios compatibles con
S3.

% RNF-15 -------------------------------------------------------------------
\item \textbf{Paginación.} \label{rnf:pagination}
Todos los \textit{endpoints} de listado de la \acs{api} deben devolver resultados paginados
mediante paginación por \textit{offset}. El tamaño de página por defecto y el tamaño máximo
son configurables. Cada respuesta paginada incluye los metadatos de paginación: número
total de elementos, página actual, tamaño de página y número total de páginas.
